% ============================================================================
% config.tex
% ----------------------------------------------------------------------------
% Este arquivo concentra as VARIÁVEIS DE CONFIGURAÇÃO VISUAL e os DADOS FIXOS
% da empresa/entidade que aparecem em todo o documento (capa, cabeçalho,
% rodapé, folha de rosto etc.).
%
% IMPORTANTE: este é o principal arquivo que você deve editar quando for
% adaptar o modelo para OUTRA empresa/marca (cores, logos, dados de contato).
% Os dados que mudam a cada documento específico (título, cliente, projeto)
% ficam em "edicao.tex", não aqui.
% ============================================================================

% -- Cor institucional -------------------------------------------------------
% Define a cor "verde-ej" usada nas linhas do cabeçalho/rodapé.
% Troque os valores RGB para adaptar à cor da sua empresa.
\definecolor{verde-ej}{RGB}{41, 128, 0}

% -- Caminhos dos arquivos de logo -------------------------------------------
% Devem apontar para arquivos de imagem existentes dentro do projeto.
% Se o arquivo não existir, o modelo mostra automaticamente um texto
% "(logo)" no lugar (ver uso de \IfFileExists em capa.tex e cabecalho-rodape.tex).
\newcommand{\logoCapaArquivo}{imag/logos/EletronLogo.jpg}                 % logo grande da capa
\newcommand{\logoCabecalhoArquivo}{imag/logos/EletronLogo_Horizontal.jpg} % logo pequena do cabeçalho

% -- Tamanhos dos logos -------------------------------------------------------
\newcommand{\larguraLogoCapa}{11cm}       % largura do logo na capa
\newcommand{\alturaLogoCabecalho}{1cm}  % altura do logo no cabeçalho

% -- Logos de parceiros/apoiadores (aparecem na capa, rodapé da página inicial)
\newcommand{\logoParceiraUmArquivo}{imag/logos/Logo_UnB.png}
\newcommand{\logoParceiraDoisArquivo}{imag/logos/Logo_Brasil_Junior.png}
\newcommand{\logoParceiraTresArquivo}{imag/logos/Logo_Concentro.png}

\newcommand{\alturaLogosParceiras}{1.5cm}     % altura padrão dos logos de parceiros
\newcommand{\opacidadeLogosParceiras}{0.75}    % opacidade (0 = invisível, 1 = opaco)

% -- Local e data -------------------------------------------------------------
\newcommand{\cidadeDocumento}{Brasília}  % cidade exibida na capa
\newcommand{\anoDocumento}{\the\year}    % ano atual, obtido automaticamente do sistema

% -- Dados fixos da empresa/entidade emissora ---------------------------------
% Usados no rodapé de todas as páginas do corpo do documento.
\newcommand{\razaoSocialEmpresa}{ELETRONJUN EMPRESA JÚNIOR DE ENGENHARIA ELETRÔNICA}
\newcommand{\cnpjEmpresa}{27.420.752/0001-68}
\newcommand{\empresaEndereco}{AE de Indústria Projeção A – UnB - Gama/DF - CEP: 72.444-240}
\newcommand{\telefoneEmpresa}{(61) 97402-7117}
\newcommand{\siteEmpresa}{www.eletronjun.com.br}

% Símbolo usado para separar os itens do rodapé (ex.: "Endereço | Telefone | Site")
\newcommand{\separadorRodape}{ \textbar\ }

% Nome completo da entidade executora do projeto (aparece na folha de rosto)
\newcommand{\entidadeExecutora}{EletronJun Empresa Júnior de Engenharia Eletrônica}

% Data de emissão: por padrão usa a data atual do sistema (\today).
\newcommand{\dataEmissaoDocumento}{\today}
